\lecture{CG(3) assembly: the Hermite stiffness matrix}{Lecture 12 --- Wednesday, September 23, 2026}{Math 6660 --- \texttt{12th.pdf}}

\noindent{\small\itshape Lecture 12 re-states the beam problem from
Lecture 11 and then carries the CG(3) discretisation all the way to the
assembled linear system. Pages 7--9 of \texttt{missing.pdf} cover the same
material; the two sources agree throughout, and the transcription below follows
my own notes.}

\begin{roadmap}[Lecture 12]
\roadq{What does the beam stiffness matrix actually look like once assembled?}
\begin{roadlist}
\item Two kinds of unknown per node (a value, a slope) means every matrix entry
  is one of four pairings. The matrix splits into blocks: $\bm{A}$
  deflection--deflection, $\bm{B}$ rotation--rotation, $\bm{C}$ the coupling.
\item $\bm{C}$ is \emph{anti}-symmetric --- and the assembled $\bm{S}$ is still
  symmetric. Worth understanding rather than memorising.
\item The load side now needs \textbf{two} vectors: a distributed load lumps
  into nodal forces \emph{and} nodal moments.
\item The numbers that come out --- $12/h^{3}$, $6/h^{2}$, $4/h$, $2/h$ --- are
  the textbook Euler--Bernoulli element matrix with $EI=1$.
\end{roadlist}
\payoff{Recognising the beam element you already know, written in the
mathematics department's notation.}
\end{roadmap}

% ==================================================================
\subsection*{Recap: the beam problem and its CG(3) space}

\begin{equation}\label{eq:l12beam}
  \left\{
  \begin{aligned}
    u^{(iv)}&=f &&\text{in }(0,1),\\
    u(0)=u(1)&=u'(0)=u'(1)=0 ,
  \end{aligned}\right.
\end{equation}
\begin{equation}\label{eq:l12V}
  \Vsp=\Bigl\{v:\ \int_0^1\bigl[v^{2}+(v')^{2}+(v'')^{2}\bigr]\dx<\infty,\ \
  v(0)=v(1)=v'(0)=v'(1)=0\Bigr\}.
\end{equation}
Functions in $\Vsp$ are $C^{1}$. Find $u\in\Vsp$ such that
\begin{equation}\label{eq:l12weak}
  \int_0^1 u''v''\dx=\int_0^1 f v\dx\qquad\forall\,v\in\Vsp .
\end{equation}

\medskip\noindent
\textbf{FEM:} $\Vh\subset\Vsp$, spanned by the two Hermite shapes per node ---
$\phih_{j}^{(1)}$, which carries the nodal value, and $\phih_{j}^{(2)}$, which
carries the nodal slope (Lecture 11, \eqref{eq:l11phi1}--\eqref{eq:l11phi2}):
\begin{equation}\label{eq:l12span}
  \Vh=\operatorname{span}\bigl\{\phih_{1}^{(1)},\dots,\phih_{M}^{(1)},\
      \phih_{1}^{(2)},\dots,\phih_{M}^{(2)}\bigr\},\qquad\dim\Vh=2M .
\end{equation}
\textbf{Trial functions:}
\begin{equation}\label{eq:l12trial}
  U=\sum_{j=1}^{M}\xi_{j}\,\phih_{j}^{(1)}(x)
    +\sum_{j=1}^{M}\eta_{j}\,\phih_{j}^{(2)}(x).
\end{equation}
\textbf{Test functions:} $v=\phih_{1}^{(1)},\dots,\phih_{M}^{(1)},
\phih_{1}^{(2)},\dots,\phih_{M}^{(2)}$.

% ==================================================================
\subsection*{The interpolant and its errors}

If $u\in\Vsp$, then
\begin{equation}\label{eq:l12PI}
  \Pi_h u=\sum_{j=1}^{M}u(x_{j})\,\phih_{j}^{(1)}(x)
          +\sum_{j=1}^{M}u'(x_{j})\,\phih_{j}^{(2)}(x),
  \qquad\text{where }\Pi_h u\in\Vh .
\end{equation}

\medskip\noindent
\textbf{Interpolation errors:}
\begin{align}
  \norm{u-\Pi_h u}_{\Ltwo[0,1]}     &\le c_{1}\norm{h^{4}u^{(4)}}_{\Ltwo[0,1]},
     \label{eq:l12int0}\\
  \norm{(u-\Pi_h u)'}_{\Ltwo[0,1]}  &\le c_{2}\norm{h^{3}u^{(4)}}_{\Ltwo[0,1]},
     \label{eq:l12int1}\\
  \norm{(u-\Pi_h u)''}_{\Ltwo[0,1]} &\le c_{3}\norm{h^{2}u^{(4)}}_{\Ltwo[0,1]}.
     \label{eq:l12int2}
\end{align}

\begin{intuition}[why $u^{(4)}$, and why the exponents step down]
The Hermite interpolant matches value and slope at both ends of every element,
so it reproduces \emph{every cubic exactly}. The first thing it cannot
reproduce is the fourth derivative --- which is why $u^{(4)}$, and nothing
lower, sits on every right-hand side. Each derivative you take of the error
costs one power of $h$: $h^{4}$ for the function, $h^{3}$ for the slope,
$h^{2}$ for the curvature. Since the energy norm for the beam measures $v''$,
\eqref{eq:l12int2} is the one that controls the FEM error, giving $O(h^{2})$.
\end{intuition}

% ==================================================================
\subsection*{Example: uniform mesh, $h_{j}=h$ for all $j$}

With
\begin{equation}\label{eq:l12fem}
  \int_0^1 U''v''\dx=\int_0^1 f v\dx,
  \qquad v=\phih_{1}^{(1)},\dots,\phih_{M}^{(1)},
            \phih_{1}^{(2)},\dots,\phih_{M}^{(2)} ,
\end{equation}
the stiffness matrix comes out of four families of integrals. Testing against
$\phih_{k}^{(1)}$, $k=1,2,\dots,M$:
\begin{equation}\label{eq:l12rowa}
  \int_0^1 U''v''\dx
  =\sum_{j=1}^{M}\xi_{j}\int_0^1\bigl(\phih_{j}^{(1)}\bigr)''
     \bigl(\phih_{k}^{(1)}\bigr)''\dx
  +\sum_{j=1}^{M}\eta_{j}\int_0^1\bigl(\phih_{j}^{(2)}\bigr)''
     \bigl(\phih_{k}^{(1)}\bigr)''\dx ,
\end{equation}
and then again, testing against $\phih_{k}^{(2)}$, $k=1,2,\dots,M$:
\begin{equation}\label{eq:l12rowb}
  \sum_{j=1}^{M}\xi_{j}\int_0^1\bigl(\phih_{j}^{(1)}\bigr)''
     \bigl(\phih_{k}^{(2)}\bigr)''\dx
  +\sum_{j=1}^{M}\eta_{j}\int_0^1\bigl(\phih_{j}^{(2)}\bigr)''
     \bigl(\phih_{k}^{(2)}\bigr)''\dx .
\end{equation}
That is a block matrix
\begin{equation}\label{eq:l12block}
  \left[\begin{array}{c|c}
    \bm{A} & \bm{C}\\\hline
    -\bm{C} & \bm{B}
  \end{array}\right]\in\R^{2M\times2M},
  \qquad\text{here }\bm{A},\bm{B},\bm{C}\in\R^{M\times M}.
\end{equation}

\begin{intuition}[where the four blocks come from]
There are two kinds of unknown ($\xi$, a value; $\eta$, a slope) and two kinds
of test function, so every entry of the matrix is one of four pairings:
value--value, value--slope, slope--value, slope--slope. $\bm{A}$ is the first,
$\bm{B}$ the last, $\bm{C}$ the cross terms. Nothing about the physics changed
--- it is still $\int u''v''$ --- only the bookkeeping, because each node now
carries two numbers.
\end{intuition}

% ------------------------------------------------------------------
\subsubsection*{The three blocks}

\medskip\noindent
Now,
\begin{equation}\label{eq:l12Ajk}
  A_{jk}=\int_0^1\bigl(\phih_{j}^{(1)}\bigr)''\bigl(\phih_{k}^{(1)}\bigr)''\dx,
  \qquad
  \bm{A}=\frac{12}{h^{3}}
  \begin{bmatrix}
    2 & -1 &        &    \\
   -1 &  2 & \ddots &    \\
      & \ddots & \ddots & -1\\
      &    & -1 & 2
  \end{bmatrix},
\end{equation}
which is $M\times M$, tridiagonal, symmetric.

\begin{equation}\label{eq:l12Bjk}
  B_{jk}=\int_0^1\bigl(\phih_{j}^{(2)}\bigr)''\bigl(\phih_{k}^{(2)}\bigr)''\dx,
  \qquad
  \bm{B}=\frac{2}{h}
  \begin{bmatrix}
    4 & 1 &        &   \\
    1 & 4 & \ddots &   \\
      & \ddots & \ddots & 1\\
      &   & 1 & 4
  \end{bmatrix},
\end{equation}
which is $M\times M$, tridiagonal, symmetric.

\begin{equation}\label{eq:l12Cjk}
  C_{jk}=\int_0^1\bigl(\phih_{j}^{(1)}\bigr)''\bigl(\phih_{k}^{(2)}\bigr)''\dx,
  \qquad
  \bm{C}=\frac{6}{h^{2}}
  \begin{bmatrix}
    0 & 1 &        &   \\
   -1 & 0 & \ddots &   \\
      & \ddots & \ddots & 1\\
      &   & -1 & 0
  \end{bmatrix},
\end{equation}
which is $M\times M$, tridiagonal --- but \emph{not} symmetric. It is
anti-symmetric:
\begin{equation}\label{eq:l12anti}
  \bm{C}^{T}=-\bm{C}.
\end{equation}
Therefore
\begin{equation}\label{eq:l12S}
  \bm{S}=\left[\begin{array}{c|c}
    \bm{A} & \bm{C}\\\hline
    \bm{C}^{T} & \bm{B}
  \end{array}\right],
  \qquad \bm{S}^{T}=\bm{S},
\end{equation}
so $\bm{S}$ is $2M\times2M$ and symmetric.

\begin{pitfall}[an anti-symmetric block does not break symmetry]
$\bm{C}^{T}=-\bm{C}$ is exactly what makes $\bm{S}$ symmetric: the lower-left
block \emph{is} $\bm{C}^{T}$, so writing it as $-\bm{C}$ in
\eqref{eq:l12block} is the same statement. The minus sign is a consequence of
the sign convention for the slope degree of freedom, not a loss of symmetry ---
$\int u''v''$ is self-adjoint, so the matrix could not have come out
unsymmetric.
\end{pitfall}

\begin{mechanics}
The $12/h^{3}$, $6/h^{2}$, $4/h$, $2/h$ pattern in
\eqref{eq:l12Ajk}--\eqref{eq:l12Cjk} is the classical Euler--Bernoulli beam
element stiffness matrix with $EI=1$. In a structural code the element matrix
is written once as
$\frac{EI}{L^{3}}\begin{bsmallmatrix}12&6L&-12&6L\\6L&4L^{2}&-6L&2L^{2}\\
-12&-6L&12&-6L\\6L&2L^{2}&-6L&4L^{2}\end{bsmallmatrix}$
and assembled; here the same numbers arrive sorted by \emph{type} of degree of
freedom instead of by node, which is why they appear as three separate blocks.
\end{mechanics}

% ==================================================================
\subsection*{RHS}

Interpolate the load with the same Hermite basis,
\begin{equation}\label{eq:l12fint}
  f\ \longrightarrow\ \widetilde{\Pi}_h f
  =\sum_{j=0}^{M+1}f(x_{j})\,\phih_{j}^{(1)}(x)
   +\sum_{j=0}^{M+1}f'(x_{j})\,\phih_{j}^{(2)}(x),
\end{equation}
\begin{equation}\label{eq:l12fsub}
  \int_0^1 f v\dx\ \longrightarrow\ \int_0^1\bigl(\widetilde{\Pi}_h f\bigr)v\dx .
\end{equation}
Now, for $j=1,2,\dots,M$,
\begin{align}
  f_{j}^{(1)}&=\frac{h}{70}\Bigl[9f(x_{j-1})+52f(x_{j})+9f(x_{j+1})\Bigr]
     +\frac{13h^{2}}{420}\Bigl[f'(x_{j-1})-f'(x_{j+1})\Bigr],
     \label{eq:l12f1}\\[0.6ex]
  f_{j}^{(2)}&=\frac{13h^{2}}{420}\Bigl[f(x_{j+1})-f(x_{j-1})\Bigr]
     -\frac{h^{3}}{420}\Bigl[3f'(x_{j-1})-8f'(x_{j})+3f'(x_{j+1})\Bigr].
     \label{eq:l12f2}
\end{align}

% ==================================================================
\subsection*{The final system}

\begin{equation}\label{eq:l12system}
  \left[\begin{array}{c|c}
    \bm{A} & \bm{C}\\\hline
    -\bm{C} & \bm{B}
  \end{array}\right]
  \left[\begin{array}{c}\bm{\xi}\\\hline\bm{\eta}\end{array}\right]
  =\left[\begin{array}{c}\bm{f}^{(1)}\\\hline\bm{f}^{(2)}\end{array}\right],
\end{equation}
which is $2M\times2M$, symmetric, positive definite. Here $\bm{\eta}$ is the
vector of \emph{derivative values at the nodes}.

\begin{mechanics}
\eqref{eq:l12system} is the assembled beam problem in the form a structural
solver produces it: $\bm{\xi}$ are the nodal deflections, $\bm{\eta}$ the nodal
rotations, $\bm{f}^{(1)}$ the consistent nodal forces and $\bm{f}^{(2)}$ the
consistent nodal moments. The $13h^{2}/420$ and $h^{3}/420$ coefficients in
\eqref{eq:l12f1}--\eqref{eq:l12f2} are the consistent load vector for a
distributed load on a Hermite beam element --- the same numbers a code computes
when you apply a pressure rather than a point load.
\end{mechanics}

\begin{dictionary}[beam FEM]
$\xi_{j}$, the coefficient of $\phih_{j}^{(1)}$ & the nodal deflection $w_{j}$\\
$\eta_{j}$, the coefficient of $\phih_{j}^{(2)}$ & the nodal rotation $\theta_{j}$\\
$\bm{A}$ block & deflection--deflection stiffness\\
$\bm{B}$ block & rotation--rotation stiffness\\
$\bm{C}$ block & the deflection--rotation coupling\\
$\bm{S}$ symmetric positive definite & the structure is stable and the energy is a genuine minimum\\
$\widetilde{\Pi}_h f$ & lumping the distributed load onto nodal forces and moments\\
\end{dictionary}

\todoq{On page~2 the first interpolation error is written
$\norm{u-\Pi_h u}\le c_{1}\norm{h^{4}u^{4}}$ --- the exponent on $u$ should be
the superscript $(4)$, matching the other two lines; transcribed as $u^{(4)}$.
Similarly the second sum in $\Pi_h u$ on that page runs to $n$ rather than $M$.}
